\subsection{Discrete-Time SIR Dynamics}

The compartments evolve according to the discrete-time SIR model:
\begin{align}\begin{split}\label{eq:sir}
S_i(t+1) &= S_i(t) - \beta(t) S_i(t) \sum_j C_{ij} I_j(t) \\
I_i(t+1) &= I_i(t) + \beta(t) S_i(t) \sum_j C_{ij} I_j(t) - \gamma I_i(t) \\
R_i(t+1) &= R_i(t) + \gamma I_i(t)
\end{split}\end{align}

We consider two regions with connectivity matrix:
\begin{equation}
\mathbf{C} = \begin{bmatrix} 1-\theta & \theta \\ \theta & 1-\theta \end{bmatrix}
\end{equation}

%% The dynamics become:
%% \begin{align}
%% S_1(t+1) &= S_1(t) - \beta(t) S_1(t)[(1-\theta) I_1(t) + \theta I_2(t)] \\
%% S_2(t+1) &= S_2(t) - \beta(t) S_2(t)[\theta I_1(t) + (1-\theta) I_2(t)] \\
%% I_1(t+1) &= I_1(t) + \beta(t) S_1(t)[(1-\theta) I_1(t) + \theta I_2(t)] - \gamma I_1(t) \\
%% I_2(t+1) &= I_2(t) + \beta(t) S_2(t)[\theta I_1(t) + (1-\theta) I_2(t)] - \gamma I_2(t)
%% \end{align}

Initial conditions are $S_i(0), I_i(0)$ and $R_i(0) = 1 - S_i(0) - I_i(0)$.

We consider two regions observed over $T$ discrete time steps ($t = 0,
1, \ldots, T-1$). The parameter vector to be estimated is:
\begin{equation}
\boldsymbol{\phi} = [\theta, S_1(0), I_1(0), S_2(0), I_2(0)]^T \in \mathbb{R}^5
\end{equation}
where $\theta$ is the connectivity parameter and $S_i(0), I_i(0)$ are initial conditions.

We focus our analysis on seasonal influenza, which provides extensive
surveillance data for validation. The discrete-time model requires
careful calibration of the recovery rate parameter $\gamma$.


From Keeling \& Rohani~\cite{keeling2008modeling}, the continuous-time
recovery rate for influenza is $\gamma_{\text{cont}} = 1/2.2 \text{
  days}^{-1}$. For weekly discrete-time observations, we must convert
this parameter appropriately.

Consider a scenario where no new infections occur (all susceptibles
removed). After one week:
\begin{align}
\text{Discrete model:} \quad I(t+1) &= (1-\gamma) I(t) \\
\text{Continuous model:} \quad I(t+7) &= \exp(-7\gamma_{\text{cont}}) I(t)
\end{align}

Matching the fraction of infecteds remaining after one week (7 days):
\begin{equation}
1 - \gamma = \exp\left(-\frac{7}{2.2}\right) = \exp(-3.18) \approx 0.041
\end{equation}

Therefore, the calibrated discrete recovery rate is:
\begin{equation}
\gamma = 1 - \exp\left(-\frac{7}{2.2}\right) \approx 0.959
\end{equation}


Here $\gamma$ takes the role of probability of recovery in one
week. Average disease duration is $1/ \gamma$ then (a geometric random
variable). The basic relation of reproduction number is $R_0
\frac{\beta}{\gamma}$. Hence
\begin{equation}
  R_0 = \frac{\beta(t)}{\gamma} = \frac{\beta_0 (1+\epsilon)}{\gamma}
\end{equation}
Taking $\epsilon = 0.3$ and assuming the $R_0$ is cited for peak infectiousness we find
$$
\beta_0 = \frac{R_0 \gamma}{1+\epsilon} = \frac{3.65 (1-\exp\left (-\frac{7}{2.2} \right))}{1+0.3} \approx 2.7
$$













